\section{Dijkstra's Algorithm}
\label{sec:graph-dijkstra}

\problemstatement{Given a weighted graph with \textbf{non-negative}
edge weights and a source node, find the shortest distance from the
source to every other node.}

\begin{patternbox}
\emph{``Shortest path, non-negative weights''} is solved by always
expanding the \textbf{unfinished} node with the currently
\textbf{smallest known} distance -- a greedy choice, made efficient by
storing candidate $(\textit{distance}, \textit{node})$ pairs in a
\textbf{min-heap} (priority queue) so the smallest is always available
in $O(\log V)$.
\end{patternbox}

\subsection*{Why the greedy choice is correct (and needs non-negative weights)}

When the node $u$ with the smallest tentative distance is popped, that
distance is guaranteed \textbf{final} -- no future relaxation through
any other, not-yet-finalized node $x$ could ever produce a shorter path
to $u$, because every such path would have to pass through some node
with tentative distance $\ge \textit{dist}[u]$ \emph{and then add a
non-negative edge weight}, which can only make the total $\ge
\textit{dist}[u]$. This argument breaks immediately if a negative edge
weight is allowed: a long detour through a very negative edge could
undercut the greedily-chosen shortest path, which is exactly why
Dijkstra's algorithm requires non-negative weights and Bellman-Ford
(Section~\ref{sec:graph-bellman-ford}) exists for the negative case.

\subsection*{Algorithm}

\begin{enumerate}
  \item Initialize $\textit{dist}[\textit{src}]=0$, all others
        $\infty$. Push $(0, \textit{src})$ into a min-heap.
  \item While the heap is non-empty: pop the smallest
        $(\textit{d}, \textit{node})$; if $\textit{d} >
        \textit{dist}[\textit{node}]$, this is a stale entry (a better
        path was already found), skip it.
  \item Otherwise, relax every edge out of $\textit{node}$: if
        $\textit{dist}[\textit{node}] + w < \textit{dist}[\textit{neighbor}]$,
        update and push $(\textit{dist}[\textit{neighbor}],
        \textit{neighbor})$.
\end{enumerate}

\begin{ideabox}[Why Stale Entries Are Allowed Rather Than Removed]
A node may be pushed onto the heap multiple times, once per
improvement to its tentative distance, rather than removing and
re-inserting it (which most heap implementations don't support
efficiently). The $\textit{d} > \textit{dist}[\textit{node}]$ check
simply ignores outdated copies when they're eventually popped --
correct, and simpler than maintaining a decrease-key operation.
\end{ideabox}

\subsection*{C++ implementation}

\begin{lstlisting}
#include <bits/stdc++.h>
using namespace std;

vector<int> dijkstra(int n, vector<vector<pair<int,int>>>& adj, int src) {
    vector<int> dist(n, INT_MAX);
    dist[src] = 0;
    priority_queue<pair<int,int>, vector<pair<int,int>>, greater<>> pq;
    pq.push({0, src});

    while (!pq.empty()) {
        auto [d, node] = pq.top();
        pq.pop();
        if (d > dist[node]) continue; // stale entry

        for (auto& [neighbor, weight] : adj[node]) {
            if (dist[node] + weight < dist[neighbor]) {
                dist[neighbor] = dist[node] + weight;
                pq.push({dist[neighbor], neighbor});
            }
        }
    }
    return dist;
}

int main() {
    int n = 5;
    vector<vector<pair<int,int>>> adj(n); // {neighbor, weight}
    auto addEdge = [&](int u, int v, int w) {
        adj[u].push_back({v, w});
        adj[v].push_back({u, w});
    };
    addEdge(0, 1, 4); addEdge(0, 2, 1); addEdge(2, 1, 2);
    addEdge(1, 3, 1); addEdge(2, 3, 5); addEdge(3, 4, 3);

    vector<int> dist = dijkstra(n, adj, 0);
    for (int d : dist) cout << d << " ";
    cout << endl; // 0 3 1 4 7
    return 0;
}
\end{lstlisting}

\subsection*{Dry run}

\dryrun{Take the graph above, source $0$.}

Heap: $[(0,0)]$. Pop $(0,0)$: relax $1$ ($\infty\to4$), $2$
($\infty\to1$); push both. Heap: $[(1,2),(4,1)]$. Pop $(1,2)$
(smallest): relax $1$ via $2$ ($4\to1+2=3$, improved); relax $3$ via
$2$ ($\infty\to1+5=6$); push $(3,1)$ and $(6,3)$. Heap:
$[(3,1),(4,1)_{\text{stale}},(6,3)]$. Pop $(3,1)$: relax $3$ via $1$
($6\to3+1=4$, improved); push $(4,3)$. Pop $(4,1)_{\text{stale}}$:
$4 > \textit{dist}[1]=3$, skip. Pop $(4,3)$: relax $4$ via $3$
($\infty\to4+3=7$); push. Pop $(6,3)_{\text{stale}}$: $6 >
\textit{dist}[3]=4$, skip. Pop $(7,4)$: no further relaxation. Final:
$\textit{dist}=[0,3,1,4,7]$.

\begin{complexitybox}
\textbf{Time Complexity.} $O((V+E)\log V)$ -- each edge can trigger one
heap push, and each heap operation costs $O(\log V)$.
\textbf{Space Complexity.} $O(V)$ for \textit{dist}, plus $O(E)$ for
the heap in the worst case (allowing stale duplicate entries).
\end{complexitybox}

\begin{takeawaybox}
Dijkstra's algorithm is BFS's natural generalization to weighted
graphs: BFS implicitly always expands the nearest unvisited node
because its queue processes things in arrival order, which only equals
distance order when every edge costs the same; Dijkstra's makes that
``always expand the nearest'' rule explicit via a priority queue,
restoring correctness for arbitrary non-negative weights.
\end{takeawaybox}
